\section{Introduction}
\label{sec:intro}

3D reconstruction and novel view synthesis are core capabilities for applications such as virtual reality, robotic navigation, and digital content creation. Neural Radiance Fields (NeRF)~\cite{mildenhall2020nerf} and its follow-ups have achieved remarkable progress in novel view synthesis and high-fidelity 3D scene reconstruction, albeit at the cost of expensive per-scene optimization. More recently, 3D Gaussian Splatting (3DGS)~\cite{kerbl3Dgaussians} represents scenes as a set of anisotropic 3D Gaussians and achieves real-time rendering with high visual quality. Building on this idea, feed-forward frameworks such as MVSplat~\cite{chen2024mvsplat} directly predict 3D Gaussian primitives from sparse multi-view images without per-scene optimization, enabling efficient reconstruction from casually captured videos.

However, most of these methods implicitly assume that the input images are \emph{clean}: captured under well-controlled conditions or carefully preprocessed, with negligible noise or degradation. In contrast, real-world multi-view data---for example, clips from web videos or mobile devices---often suffer from sensor noise, compression artifacts, and low-light degradation. Noise not only corrupts geometry estimation but also damages textures and fine details, especially when fusing information across sparse views. Under this realistic setting, 3D reconstruction from noisy multi-view inputs remains largely under-explored in the current NeRF/3DGS literature.

A straightforward engineering solution is a two-stage pipeline. One first denoises each frame in the 2D image domain using modern learning-based denoisers, such as the recent generalizable IDF network~\cite{kim2025idf} or classical CNN-based models like DnCNN~\cite{zhang2017dncnn} and FFDNet~\cite{zhang2018ffdnet}, and then feeds the denoised images into NeRF- or 3DGS-based pipelines for 3D reconstruction. While simple and modular, this design has notable drawbacks: (1) denoising can easily introduce over-smoothing and remove subtle details; (2) independently denoising each view weakens multi-view consistency, which is crucial for reliable 3D fusion; and (3) decomposing the pipeline into multiple modules increases inference latency and system complexity. At the other extreme, directly feeding noisy images into vanilla 3DGS or MVSplat leads to rapid degradation of reconstruction quality as noise intensity increases, as we will empirically demonstrate in Sec.~\ref{sec:experiments}.

In this work, we instead tackle noise \emph{inside the 3D representation}. Rather than treating denoising as an external pre-processing step, we aim to let a feed-forward 3D Gaussian splatting network itself learn to recover clean 3D scenes from noisy multi-view observations. Concretely, we build on the multi-view encoding and 3D decoding architecture of MVSplat and propose \emph{DenoiseSplat}, a feed-forward 3D Gaussian splatting method tailored to noisy inputs. Given noisy images and camera parameters as input and clean renderings as the only supervision, DenoiseSplat is trained end-to-end to predict a clean 3D Gaussian representation and its rendered views, without any access to ground-truth 3D supervision.

Fig.~1 offers a teaser overview of our works. To support this setting, we construct a multi-noise, scene-consistent noisy--clean dataset based on RealEstate10K (RE10K). For each scene, we inject one of four noise types---Gaussian, Poisson, speckle, or salt-and-pepper---into the 2D RGB images and keep the noise type and level consistent across all views within that scene, yielding paired noisy--clean multi-view samples. This simple yet controlled degradation pipeline mimics realistic acquisition conditions with scene-level noise characteristics and provides a reproducible benchmark for studying noisy 3D reconstruction.

To further improve robustness under heavy noise, we introduce a \emph{dual-branch Gaussian head} for geometry--appearance decoupling, as illustrated in Fig.~2. Instead of using a single shared head to regress all Gaussian parameters, we decouple geometry-related and appearance-related predictions into two lightweight branches: a geometric branch that outputs centers, rotations, scales, and opacities, and an appearance branch that predicts spherical-harmonics coefficients and colors. This design allows geometry to be estimated from more stable, noise-robust cues, while appearance is refined separately to absorb residual noise and color fluctuations. At test time, given only noisy multi-view inputs, a single forward pass through DenoiseSplat produces a denoised 3DGS scene and high-quality novel-view renderings, without any test-time optimization.

Our main contributions are summarized as follows:
\begin{itemize}[leftmargin=*,
                topsep=0pt,itemsep=0pt,parsep=0pt,partopsep=0pt]
    \item \textbf{Problem formulation and feed-forward framework for noisy multi-view reconstruction.}
    We systematically study 3D reconstruction from noisy multi-view inputs and, building on MVSplat, propose DenoiseSplat, a feed-forward 3D Gaussian splatting framework specifically designed for noise-corrupted images. Our method preserves the advantages of no per-scene optimization and efficient rendering, while explicitly enhancing robustness to noise through architectural and training choices.
    
    \item \textbf{Dual-branch Gaussian head for geometry--appearance decoupling.}
    We redesign the Gaussian head into a dual-branch structure, where a geometric branch predicts position-, rotation-, scale-, and opacity-related parameters and an appearance branch predicts spherical-harmonics coefficients and colors. This geometry--appearance decoupling mitigates the interference between noisy appearance and stable geometry, leading to more consistent 3D structure and sharper textures under strong noise.
    
    \item \textbf{Multi-noise, scene-consistent noisy--clean data construction.}
    We develop a simple yet reproducible noise degradation pipeline on RE10K, introducing four noise types (Gaussian, Poisson, speckle, and salt-and-pepper) and adopting a scene-level noise configuration, where all views of a scene share the same noise type and level. This setting better reflects realistic acquisition conditions and provides a useful benchmark for studying noisy 3D reconstruction.
    
    \item \textbf{Comprehensive experiments and ablations on noise levels and design choices.}
    On the constructed noisy RE10K benchmark, we compare DenoiseSplat against vanilla MVSplat and strong two-stage baselines that combine state-of-the-art 2D denoisers with 3D reconstruction under multiple noise types and intensities. We further conduct ablation studies on the standard deviation of Gaussian noise, the hyperparameters of other noise types, and the proposed dual-branch Gaussian head, together with qualitative analyses under different noise settings. Across a wide range of noise regimes, our method consistently outperforms baselines in PSNR, SSIM, and LPIPS, while preserving more texture details and geometric consistency.
\end{itemize}
